\documentclass[10pt,a4paper]{amsart}
\usepackage[margin=19mm]{geometry}
\usepackage{amsmath,amssymb,mathtools}
\usepackage[scr=rsfs,cal=euler]{mathalfa}
\usepackage{xfrac}
\usepackage[unicode,hidelinks,bookmarks=false]{hyperref}
\newcommand{\ie}{i.e.\ }
\newcommand{\cf}{cf.\ }
\newcommand{\resp}{respectively\ }
\newcommand{\Subsection}{\subsection}
\providecommand{\nobreakdash}{\nobreak\hbox{-}}
\setlength{\emergencystretch}{2em}
\multlinegap0pt
\usepackage{bm}
\usepackage{bbm}
\usepackage{dsfont}
\usepackage{xspace}
\usepackage{cancel}
\usepackage{mathtools}
\usepackage{amsthm}
\makeatletter
\@ifundefined{proposition}{\newtheorem{proposition}{Proposition}}{}
\makeatother
\newcommand{\C}{\mathbb{C}}
\newcommand{\R}{\mathbb{R}}
\title[Multiversion of the Hausdorff--Young inequality]{Multiversion of the Hausdorff--Young inequality}
\author{Paata Ivanisvili, Pavlos Kalantzopoulos}
\date{Source: arXiv:2506.08494v1 (CC BY 4.0)}
\begin{document}
\maketitle

\noindent\emph{Source and licence.} Multiversion of the Hausdorff--Young inequality. Version arXiv:2506.08494v1 (CC BY 4.0), distributed under the Creative Commons Attribution 4.0 International licence, as recorded in the arXiv licence metadata for this exact version and shown on the abstract page https://arxiv.org/abs/2506.08494v1. Full rights notice: licenses/arxiv-2506.08494-CC-BY-4.0.txt.\par\smallskip

\noindent\textit{Editorial scope.} Counted unit: the Proposition whose source label is \texttt{equivalComplex} in the article (source0.tex lines 1256--1293, source label \texttt{equivalComplex}), delivered together with the article's own complete original proof of it (source0.tex lines 1296--1595, one \texttt{proof} environment of 300 lines). No mathematics is written, repaired or re-derived: statement and proof are the article's own text line for line. Required context retained uncounted: none. Every definition, display and label of those lines is retained unchanged. Omitted from this package and not redistributed: 1--1255, 1294--1295, 1596--2278. Nothing inside a retained excerpt is omitted or altered: every retained body is delivered exactly as the article's lines are written, so no omission receipt of any kind accompanies this package. Citations: the retained excerpts carry no citation command at all, so no bibliography entry of the article is transcribed and no reference is redistributed. Reference adaptation: none was needed and none was made; every reference inside the delivered unit resolves inside it, so no unresolved reference or citation is produced. Printed slips kept literal: no printed word, symbol, hypothesis or number of the counted statement or of its proof was altered. Layout and class: the article's own class is \texttt{amsart} with packages including amsfonts, amssymb, amsmath, amsthm, color, hyperref, pdfsync, comment, calrsfs; the wrapper is the lane's pinned \texttt{amsart} editorial wrapper at 10pt on A4, which restates the theorem-like environments the retained text needs and adds the article's own macro definitions that the retained text needs (\texttt{\textbackslash C}, \texttt{\textbackslash R}), with extra wrapper packages (bm, bbm, dsfont, xspace, cancel, mathtools). The delivered numbering is the unit's own. Assets: no retained span contains a figure environment, a graphics-inclusion command, a PDF-page inclusion, a table or a TikZ picture; no image file is copied into this package, the package declares no asset dependency and no image file is redistributed with it. Licence: version arXiv:2506.08494v1 of this article is distributed under the Creative Commons Attribution 4.0 International licence, as recorded in the arXiv licence metadata for this exact version and shown on the abstract page https://arxiv.org/abs/2506.08494v1. A full rights notice is delivered at licenses/arxiv-2506.08494-CC-BY-4.0.txt. Characters: every retained excerpt is pure ASCII, so the delivered engine, XeLaTeX with its default Latin Modern text font, reports no missing character.
\begin{proposition}\label{equivalComplex}
Let $F :[0,\infty) \mapsto \mathbb{R}$ and $M: [0,\infty)^{n} \mapsto [0,\infty)$ be two functions that are smooth in the interior of their domains, continuous up to the boundary,  and satisfy polynomial growth conditions. We also assume that $F'>0$ on $(0,\infty)$, the map 
\begin{align}\label{lemmaConvF}
(t,y) \mapsto \frac{F''(t)}{F'(t)}y^{2}
\end{align}
is convex on $(0,\infty)\times \mathbb{R}$, and 
 $M_{m}>0$ on $(0, \infty)^{n}$ for all $m=1, \ldots, n$.    Let $z=(z_{1}, \ldots, z_{n}) \in \mathbb{C}^{n}$ with $|z_{j}|\leq 1$, $j=1, \ldots, n$. The following are equivalent.
\vspace{0.2cm}
\begin{enumerate}
\item \label{Globalcomplex} \textbf{Global:} For all polynomials $f_j:\R^{k_j}\to\C$, $j=1, \ldots, n$,  it holds
\vspace{0.2cm}
\begin{align*}
\int_{\R^k}F\circ M(|T_{z_1}&(f_1(A_1 x))|^2, \ldots, |T_{z_n}(f_n( A_n x))|^2) \,d\gamma(x)\\
&\leq F\left(	\int_{\R^k}M(|f_1(A_1 x)|^2, \ldots, |f_n(A_n x)|^2)\,d\gamma(x)\right).
\end{align*}  
\vspace{0.2cm}
\item \label{Localcomplex}\textbf{Local:}  For all ${\bf{c}}=(c_{1},\ldots, c_{n})\in (0,\infty)^n$ it holds
\vspace{0.2cm}
\begin{align*}
		&-\frac{F''(M({\bf{c}}))}{F'(M({\bf{c}}))}\left(\sum_{m, q=1}^{n} M_{m}({\bf{c}})M_{q}({\bf{c}})  (\Re(z_mw_{m}))^T A_{m} A_{q}^T  \Re(z_qw_{q}) \right)\\
		&+\left(\sum_{m, q=1}^{n}  M_{mq}({\bf{c}}) \left( (\Re(w_{m}))^T A_{m} A_{q}^T  \Re(w_{q})  -  (\Re(z_mw_{m}))^T A_{m} A_{q}^T  \Re(z_qw_{q}) \right) \right)\\
		&+\left(\frac{1}{2}\sum_{m=1}^n\frac{M_{m}({\bf{c}})}{c_m}(1-|z_m|^2) 
		|w_{m}|^2\right)\geq0
	\end{align*}
	for all $w_{m}\in\C^{k_m}$, $1\leq m\leq n$
\vspace{0.2cm}
\item \label{monotonicitycomplex} \textbf{Monotonicity:} Let $\ell_p:\R^{k_p}\to\mathbb{C}$, $p=1,\ldots,n$ be any polynomials. Consider the function 
\begin{equation}\label{gforComplex}
g_p(u,x,s)=Q_{1-s}^yQ_s^v\ell_p((A_p u+iv)+z_p(A_p x+iy)),
\end{equation}
where $u,x\in\R^k$, $s\in[0,1]$. Then the function $C:[0,1]\to \R$ given by 
\begin{align}\label{4Hflows}
C(s)=Q_{1-s}^x F\big(Q_s^uM\big(|g_1(u,x,s)|^2,\ldots,|g_n(u,x,s)|^2)\big)
\end{align}
is nondecreasing. 
\vspace{0.2cm}
\end{enumerate} 
\end{proposition}	

\begin{proof}
\eqref{Globalcomplex}$\implies$\eqref{Localcomplex}:		Consider  $f_j:\R^{k_j}\to\C$, $j=1,\ldots,n$, 
		\begin{equation*}
			f_j(x_j) = a_j+\varepsilon (\eta_j\cdot x_j),
		\end{equation*}
		where $a_j\in \mathbb{C}$ are nonzero, $\eta_j\in\C^{k_j}$ and  $\varepsilon >0$. Notice that   $T_{z_j}f_j:\R^{k_j}\to\C$ are given by
		\begin{equation*}
			T_{z_j} f_j(x_j)= a_j+\varepsilon z_j(\eta_j\cdot x_j).
		\end{equation*}
		Set $J(x):=F(M(x))$, $x\in(0,\infty)^n$ and also $g_j(x):=T_{z_j} f_j(A_jx)$, $x\in\R^k$. By the growth condition on $F$ and $M$ we have  that
		\begin{equation*}
			\int_{\R^n}J(|g_1|^2,\ldots,|g_n|^2)^2\,d\gamma(x)\leq C,
		\end{equation*}
		and thus by Cauchy--Schwarz inequality we see that
		\begin{align*}
			\int_{| x|> \varepsilon^{-\frac{1}{2}}}J(|g_1|^{2},\ldots,|g_n|^2)\,d\gamma(x)=\mathcal{O}(e^{-\frac{1}{100\varepsilon}}),
		\end{align*}
		and hence 
		\begin{align*}
			\int_{\R^k}J(|g_1|^{2},\ldots,|g_n|^2)\,d\gamma(x)=\int_{| x|\leq \varepsilon^{-\frac{1}{2}}}J(|g_1|^{2},\ldots,|g_n|^2)\,d\gamma(x)+\mathcal{O}(\varepsilon^3).
		\end{align*}
		For $a\in (0,\infty)^{n}$, Taylor's formula implies that
		\begin{equation*}
			J(a+\delta) = J(a) + \sum_{j=1}^n J_j(a) \delta_j + \frac{1}{2} \sum_{ i,j=1}^{n} J_{ij}(a) \delta_i\delta_j + \mathcal{O}(\sum_{j=1}^n|\delta_j|^3),   
		\end{equation*}
		provided that the norm of $\delta=(\delta_1,\ldots,\delta_n)\in\R^n$ is small enough. We write
		\begin{align*}
			|g_j(x)|^2=|a_j|^2+2\varepsilon \Re(\bar{a}_jz_j(\eta_j\cdot A_jx))+\varepsilon^2|z_j(\eta_j\cdot A_j x)|^2
		\end{align*}
		and we see that as  $\varepsilon\to0$, for any $|x|\leq \varepsilon^{-1/2}$ the norm of the vector
		\begin{equation*}
			\{\delta_j(x)\}_{j=1}^n:=\{2\varepsilon \Re(\bar{a}_jz_j(\eta_j\cdot A_jx))+\varepsilon^2|z_j(\eta_j\cdot A_j x)|^2\}_{j=1}^n
		\end{equation*} 
		goes to zero,  and so for small $\varepsilon>0$   we have 
        
		\begin{align*}
			J(|g_1|^2,\ldots,|g_n|^2)&= J(|a_1|^2,\ldots,|a_n|^2)  +\sum_{j=1}^n J_j(|a_1|^2,\ldots,|a_n|^2) \delta_j(x) \\
			&+ \frac{1}{2} \sum_{ i,j=1}^{n} J_{ij}(|a_1|^2,\ldots,|a_n|^2) \delta_i(x)\delta_j(x) + \mathcal{O}(\sum_{j=1}^n|\delta_j(x)|^3),   
		\end{align*}
		provided that $|x|\leq \varepsilon^{-1/2}$. Using the identity 
		\begin{align*}
			\int_{\R^n}(\eta_i\cdot A_ix)( \eta_j\cdot A_jx)\,d\gamma(x)=\eta_i^TA_iA_j^T \eta_j,
		\end{align*}
		and the fact that $\delta_j$ is a polynomial, we obtain that
		\begin{align*}
			\int_{| x|\leq \varepsilon^{-\frac{1}{2}}}\delta_j(x)\,d\gamma(x)&=\int_{\R^n}2\varepsilon \Re(\bar{a}_jz_j(\eta_j\cdot A_jx))+\varepsilon^2|z_j(\eta_j\cdot A_j x)|^2\,d\gamma(x)+ \mathcal{O}(\varepsilon^3)\\
			&=\varepsilon^2 |z_j|^2\|\eta_j\|^2 + \mathcal{O}(\varepsilon^3).
		\end{align*}
		Similarly, for small $\varepsilon>0$ we have that
		\begin{align*}
			\int_{| x|\leq \varepsilon^{-\frac{1}{2}}}\delta_i\delta_j\,d\gamma(x)
			&=  4\varepsilon^2\int_{\R^k}\Re(\bar{a}_iz_i(\eta_i\cdot A_ix)) \Re(\bar{a}_jz_j(\eta_j\cdot A_jx))\,d\gamma(x)+\mathcal{O}(\varepsilon^3)\\
			&= 4\varepsilon^{2} (\Re(\bar{a}_{i}z_{i}\eta_{i}))^{T} A_{i}A_{j}^{T} \Re(\bar{a}_{j}z_{j}\eta_{j}) + \mathcal{O}(\varepsilon^3).
		\end{align*} 
		Denote $c=(|a_1|^2,\ldots,|a_n|^2)$. The previous two calculations imply
		\begin{align*}
			\int_{| x|\leq \varepsilon^{-\frac{1}{2}}}J(|g_1|^2,\ldots,|g_n|^2)\,d\gamma(x)&= J(c)  +\varepsilon^2\sum_{j=1}^n J_j(c) |z_j|^2\|\eta_j\|^2\\
			&+ 2\varepsilon^2\sum_{1\leq i,j\leq n} J_{ij}(c)(\Re(\bar{a}_{i}z_{i}\eta_{i}))^{T} A_{i}A_{j}^{T} \Re(\bar{a}_{j}z_{j}\eta_{j})\\
   &+ \mathcal{O}(\varepsilon^3).
		\end{align*}
		Let $q_j(x)=f_j(A_jx)$, $x\in\R^k$, $j=1,\ldots,n$. Similarly,	since
		\begin{align*}
			|q_j(x_j)|^2&=|a_j|^2+2\varepsilon \Re(\bar{a}_j(\eta_j\cdot A_jx_j))+\varepsilon^2|\eta_j\cdot A_jx_j|^2
		\end{align*}
		we obtain
		\begin{align*}
			\int_{| x|\leq \varepsilon^{-\frac{1}{2}}}M(|q_1|^2,\ldots,|q_n|^2)d\gamma &= M(c)  +\varepsilon^2\sum_{j=1}^n M_j(c) \|\eta_j\|^2\\
			&+ 2\varepsilon^2\sum_{1\leq i,j\leq n} M_{ij}(c)(\Re(\bar{a}_{i}\eta_{i}))^{T} A_iA_j^T\Re(\bar{a}_{j}\eta_{j})\\
   &+ \mathcal{O}(\varepsilon^3).
		\end{align*}
		For $t_{0}=M(c)$ we have 
		\begin{align*}
			F(t_{0}+\varepsilon_{1}) = F(t_{0}) +F'(t_{0})\varepsilon_{1}+\mathcal{O}(\varepsilon_{1}^{2}),
		\end{align*}
		as $\varepsilon_{1} \to 0^{+}$. Thus, we have
		\begin{align*}
			F\left( \int_{\mathbb{R}} M(|f_1|^2,\ldots,|f_n|^2) d\gamma\right) &  = F(M(c))\\
			&\quad+F'(M(|a|^2))\varepsilon^2\sum_{j=1}^n M_j(c) \|\eta_j\|^2\\
			&\quad +F'(M(c)) 2\varepsilon^2\sum_{1\leq i,j\leq n} M_{ij}(c)(\Re(\bar{a}_{i}\eta_{i}))^{T} A_iA_j^T \Re(\bar{a}_{j}\eta_{j})  + \mathcal{O}(\varepsilon^3).
		\end{align*}
		The complex global inequality implies that
		\begin{align*}
			& J(c)  +\varepsilon^2\sum_{j=1}^n J_j(c) |z_j|^2\|\eta_j\|^2\\
			&+ 2\varepsilon^2\sum_{ i,j=1}^{n} J_{ij}(c)(\Re(\bar{a}_{i}z_{i}\eta_{i}))^{T} A_{i}A_{j}^{T} \Re(\bar{a}_{j}z_{j}\eta_{j}) + \mathcal{O}(\varepsilon^3)\\
			&\leq F(M(c))+F'(M(c))\varepsilon^2\sum_{j=1}^n M_j(c) \|\eta_j\|^2\\
			& +F'(M(c)) 2\varepsilon^2 \sum_{ i,j=1}^{n} M_{ij}(c)(\Re(\bar{a}_{i}\eta_{i}))^{T} A_iA_j^T\Re(\bar{a}_{j}\eta_{j})  + \mathcal{O}(\varepsilon^3).
		\end{align*}
		Recall, that $J(c)=F(M(c))$. So canceling the constant terms,  dividing both sides of the inequality by $\varepsilon^{2}$,  and taking $\varepsilon \to 0$ we obtain 
		\begin{align*}
			& \sum_{j=1}^n J_j(c) |z_j|^2\|\eta_j\|^2+ 2 \sum_{ i,j=1}^{n} J_{ij}(c)(\Re(\bar{a}_{i}z_{i}\eta_{i}))^{T} A_{i}A_{j}^{T} \Re(\bar{a}_{j}z_{j}\eta_{j})\leq\\
			&\leq F'(M(c))\sum_{j=1}^n M_j(c) \|\eta_j\|^2+ 2F'(M(c)) \sum_{ i,j=1}^{n} M_{ij}(c)(\Re(\bar{a}_{i}\eta_{i}))^{T} A_iA_j^T \Re(\bar{a}_{j}\eta_{j}).
		\end{align*}
		Denoting $\omega_j = \bar{a}_j\eta_j \in\C^{k_j}$, $j=1,\ldots,n$ the above rewrites
		\begin{align*}
			& \sum_{j=1}^n J_j(c) |z_j|^2\|\eta_j\|^2+ 2 \sum_{ i,j=1}^{n} J_{ij}(c)(\Re(z_{i}\omega_i))^{T} A_{i}A_{j}^{T} \Re(z_{j}\omega_j)\leq\\
			&\leq F'(M(c))\sum_{j=1}^n M_j(c) \|\eta_j\|^2+ 2 F'(M(c))\sum_{ i,j=1}^{n} M_{ij}(c)(\Re(\omega_i))^{T} A_iA_j^T \Re(\omega_j).  
		\end{align*}
		After substituting the values 
		\begin{align*}
			J_i(c) &= F'(M(c)) M_i(c),\\
			J_{ij}(c) &= F''(M(c)) M_i(c)M_j(x) + F'(M(c)) M_{ij}(c),
		\end{align*}
		and then dividing by 2$F'(M(c))$, it becomes 
		\begin{align*}
			&-\frac{F''(M(c))}{F'(M(c))} \sum_{ i,j=1}^{n} M_i(c)M_j(x) (\Re(z_{i}\omega_i))^{T} A_{i}A_{j}^{T} \Re(z_{j}\omega_j)\\
			&+\sum_{ i,j=1}^{n} M_{ij}(c)\left[ (\Re(\omega_i))^{T} A_iA_j^T\Re(\omega_j)-  (\Re(z_{i}\omega_i))^{T} A_{i}A_{j}^{T} \Re(z_{j}\omega_j) \right] \\
			&+ \frac{1}{2}\sum_{j=1}^n M_j(c)(1-|z_j|^2) \|\eta_j\|^2  \geq0
		\end{align*}
		which coincide with the local condition of the theorem.
		
\eqref{Localcomplex}$\implies$\eqref{monotonicitycomplex}: It suffices to show the implication for a translation $\tilde{M}(t_1,\ldots,t_n) = M(t_1+\varepsilon,\ldots,t_n+\varepsilon)$, where $\varepsilon > 0$ is fixed, instead of $M$. The main reason for introducing $\tilde{M}$ is to avoid issues with undefined values such as $M_m(0,\ldots,0)$, $M_{mq}(0,\ldots,0)$, and the fact that $M(0,\ldots,0)$ can be zero  which can make $F'(M(0,\ldots,0))$ and $F''(M(0,\ldots,0))$ undefined. 
		
		First we notice that the local condition for $M$ implies the local condition for $\tilde{M}$. Indeed, notice that $\tilde{M} > M(\varepsilon, \ldots, \varepsilon) > 0$ since $M$ has all partial derivatives positive. Next, it suffices to apply the local inequality at the point $\mathbf{t} + \varepsilon = (t_1 + \varepsilon, \ldots, t_n + \varepsilon)$ instead of $\mathbf{t} = (t_1, \ldots, t_n)$, and use the inequality
		\begin{align*}
			\sum_{m=1}^n\frac{M_{m}({\bf{t}}+\varepsilon)}{t_m+\varepsilon}\frac{1-|z_m|^2}{2}\sum_{b=1}^{k_m}|w_{m,b}|^2\leq 	\sum_{m=1}^n\frac{\tilde{M}_{m}({\bf{t}})}{t_m}\frac{1-|z_m|^2}{2}\sum_{b=1}^{k_m}|w_{m,b}|^2.
		\end{align*}
		Next, we notice that the statement about  monotonicity for $\tilde{M}$ implies the statement about monotonicity for $M$. Indeed, for any $s_{1}, s_{2} \in (0,1)$, $s_{1}<s_{2},$ we have $\tilde{C}(s_{1})\leq \tilde{C}(s_{2})$, where 
		\begin{align*}
			\tilde{C}(s):=Q_{1-s}^x F\big(Q_s^uM\big(\big|Q_{1-s}^yQ_s^v&\ell_1((A_1u+iv)+z_1(A_1x+iy)\big)\big|^2+\varepsilon,\ldots\\
			&\ldots,\big|Q_{1-s}^yQ_s^v\ell_n((A_n u+iv)+z_n(A_n x+iy))\big|^2+\varepsilon\big)\big)\geq0.
		\end{align*}
		Let us denote $t_m=\big|Q_{1-s}^yQ_s^v\ell_n((A_m u+iv)+z_m(A_m x+iy))\big|^2$, $m=1,\ldots,n$. Letting $\varepsilon \to 0$, by continuity, $M(t_1 + \varepsilon, \ldots, t_n + \varepsilon) \to M(t_1, \ldots, t_n)$. Since $M$  satisfies polynomial growth condition, applying the Lebesgue dominated convergence theorem we get also that $Q_s^{u}M(t_1+\varepsilon,\ldots,t_n+\varepsilon)\to Q_s^{u}M(t_1,\ldots,t_n)$. Then, by the continuity of $F$ we apply $F$ to this converges and last for the same reason apply $Q_{1-s}^x$, finally obtaining $C'(s) \geq 0$ for $s \in (0,1)$. 
        
        In what follows, for simplicity, we will continue the proof using the letter $M$ instead of $\tilde{M}$, with the additional assumption on $M$ that all the partial deriavtives are well defined on the boundary of the domain of $M$, and $M(0,0, \ldots, 0)>0$.
		
		For $m=1,\ldots,n$, let $\ell_m:\R^{k_m}\to\R$ be polynomials,  and $g_m$ the quantity given in \eqref{gforComplex}. Set
				\begin{align*}
			&{\bf{G}}=(|g_1|^2,\ldots,|g_n|^2),\\
			&{\bf{A}}=Q_{s}^{u}M({\bf{G}}).
		\end{align*}
		With this notation, the function given in \eqref{4Hflows} is written as $C(s)=Q_{1-s}^xF(Q_{s}^{u}M({\bf{G}}))$, and thus we have 
        \begin{align*}
	C'(s)&=\left(\frac{\mathrm{d}}{\mathrm{d}s}Q_{1-s}^x\right)F(Q_{s}^{u}M({\bf{G}}))\\
	&\quad+Q_{1-s}^xF'(Q_{s}^{u}M({\bf{G}}))\left(\frac{\mathrm{d}}{\mathrm{d} s}Q_s^u\right)M({\bf{G}}) \\
	&\quad+Q_{1-s}^xF'(Q_{s}^{u}M({\bf{G}})) Q_s^u\frac{\mathrm{d}}{\mathrm{d} s} M({\bf{G}}).
\end{align*}

In turn, by the heat equation we have
		\begin{equation*}
			C'(s)=\frac{1}{2}Q_{1-s}^x\left(-\Delta_xF({\bf{A}})+F'({\bf{A}})Q_s^u\Delta_uM({\bf{G}}) +2F'({\bf{A}}) Q_s^u\frac{\mathrm{d}}{\mathrm{d} s} M({\bf{G}})\right).
		\end{equation*}
		Let $\ell_{m,b}$ be the partial derivative of $\ell_m$ with respect to $x_b$ variable. Set
		\begin{equation*}
			g_{m,b}=Q_{1-s}^yQ_s^v    \ell_{m,b}((A_m u+iv)+z_m(A_m x+iy)).
		\end{equation*}
	Similarly we denote $g_{m.bb}$ as the corresponding functions associated with the second order partial derivatives of $\ell_{m,bb}$. Let $A_m^b$ be $b$'th row of the matrix $A_m$, $b=1,\ldots,k_m$,  and let $ e_1, \ldots, e_k $ be the standard orthonormal basis in $ \mathbb{R}^k$.	We see that the partial derivatives $(g_m)_{u_j}$ and $(g_m)_{x_j}$ are scalar multiples of each other, since
		\begin{align*}
			(g_m)_{u_j}&=\sum_{b=1}^{k_m}g_{m,b}(A_{m,b}\cdot e_j),\\
			(g_m)_{x_j}&=z_m \sum_{b=1}^{k_m}g_{m,b}(A_{m,b}\cdot e_j).
		\end{align*}
		Therefore,
	\begin{align*}
		(|g_m|^{2})_{u_j}&= 2\sum_{b=1}^{k_m}(A_{m,b}\cdot e_j) \Re(g_{m,b} \bar{g}_m ),  \\ 
		(|g_m|^{2})_{x_j}&=2\sum_{b=1}^{k_m}(A_{m,b}\cdot e_j) \Re(z_m g_{m,b} \bar{g}_m ).
	\end{align*}
	Moreover,
	\begin{align*}
		(g_m)_{u_ju_j}&=\sum_{b_1, b_2=1}^{k_m}g_{m,b_1b_2}(A_{m,b_1}\cdot e_j)(A_{m,b_2}\cdot e_j)\\
		(g_m)_{x_jx_j}&=z_m^2\sum_{b_1, b_2=1}^{k_m}g_{m,b_1b_2}(A_{m,b_1}\cdot e_j)(A_{m,b_2}\cdot e_j)
	\end{align*}
and thus we have
		\begin{align*}
			(|g_m|^2)_{u_ju_j}&=  \sum_{b_1, b_2=1}^{k_m}(A_{m,b_1}\cdot e_j)(A_{m,b_2}\cdot e_j)g_{m,b_1b_2}\bar{g}_m\\
			&+\sum_{b_1, b_2=1}^{k_m}(A_{m,b_1}\cdot e_j)(A_{m,b_2}\cdot e_j)\bar{g}_{m,b_1b_2}g_m  \\
			&+2|(g_m)_{u_j}|^2  \\
			(|g_m|^{2})_{x_jx_j}&=  \sum_{b_1, b_2=1}^{k_m}z_m^2(A_{m,b_1}\cdot e_j)(A_{m,b_2}\cdot e_j)g_{m,b_1b_2}\bar{g}_m\\
			&+\sum_{b_1, b_2=1}^{k_m}\bar{z}_m^2(A_{m,b_1}\cdot e_j)(A_{m,b_2}\cdot e_j)\bar{g}_{m,b_1b_2}g_m\\
			&  +2|(g_m)_{x_j}|^2.  
		\end{align*}
		Next, the product rule and the heat equation together imply
		\begin{align*}
			(g_m)_s&=\frac{\mathrm{d}}{\mathrm{d}s}Q_{1-s}^yQ_s^v\ell_m(S) \\
			&=\frac{1}{2}Q_{1-s}^yQ_s^v\left(-\sum_{b=1}^{k_m}\ell_{m,bb}(S)(iz_m)^2+\sum_{b=1}^{k_m}\ell_{m,bb}(S)(i)^2\right)\\
			&=\frac{z_m^2-1}{2}\sum_{b=1}^{k_m}g_{m,bb},
		\end{align*}
		where $S=( A_m u+iv)+z_m( A_m x+iy)$, and therefore, 
		\begin{align*}
			\frac{\mathrm{d}}{\mathrm{d}s}M({\bf{G}})&=\sum_{m=1}^nM_{m}({\bf{G}})\left(|g_m|^2\right)_s\\
			&=\sum_{m=1}^nM_{m}({\bf{G}})\left(\bar{g}_m(g_m)_s+g_m(\bar{g}_m)_s\right)\\
			&=\sum_{m=1}^nM_{m}({\bf{G}})\left(\frac{z^{2}_m-1}{2}\sum_{b=1}^{k_m}\bar{g}_mg_{m,bb}+\frac{\bar{z}^{2}_m-1}{2}\sum_{b=1}^{k_m}g_m\bar{g}_{m,bb}\right).
		\end{align*}
		To calculate $\Delta_xF({\bf{A}})$ first we otice that 
		\begin{equation*}
			\frac{\partial}{\partial x_j}F({\bf{A}})=F'({\bf{A}})	\frac{\partial}{\partial x_j}(Q_s^uM({\bf{G}}))=F'({\bf{A}})\left(Q_s^u\sum_{m=1}^nM_{m}({\bf{G}})(|g_m|^2)_{x_j}\right),
		\end{equation*}
		and so
		\begin{align*}
			\frac{\partial^2}{\partial x_j^2}F({\bf{A}})&=F''({\bf{A}})\left(Q_s^u\sum_{m=1}^nM_{m}({\bf{G}})(|g_m|^2)_{x_j}\right)^2\\
			&+F'({\bf{A}})Q_s^u\left(\sum_{m, q=1}^{n}M_{mq}({\bf{G}})(|g_m|^2)_{x_j}(|g_q|^2)_{x_j} \right)\\
			&+F'({\bf{A}})Q_s^u\left( \sum_{m=1}^nM_{m}({\bf{G}})(|g_m|^2)_{x_jx_j}\right).
		\end{align*}
		On the other hand, for $\Delta_uM({\bf{G}})$ we have
		\begin{align*}
			\frac{\partial^2}{\partial u_j^2}M({\bf{G}})&=\frac{\partial}{\partial u_j}\left(\sum_{m=1}^nM_{m}({\bf{G}})(|g_m|^2)_{u_j}\right)\\
			&=\sum_{m, q=1}^{n}M_{mq}({\bf{G}})(|g_m|^2)_{u_j}(|g_q|^2)_{u_j}+\sum_{m=1}^nM_{m}({\bf{G}})(|g_m|^2)_{u_ju_j}.
		\end{align*}
		We can write $C'(s)=\frac{1}{2}Q_{1-s}^x\left(\sum_{j=1}^k\mathcal{L}_j\right)$ where 
		\begin{align*}
			\mathcal{L}_j=-	\frac{\partial^2}{\partial x_j^2}F({\bf{A}})+  F'({\bf{A}}) Q_s^u \frac{\partial^2}{\partial u_j^2}M({\bf{G}})+\frac{2}{k}F'({\bf{A}})Q_s^u\frac{\mathrm{d}}{\mathrm{d}s}M({\bf{G}}).
		\end{align*}
		Substituting we see that
		\begin{align*}
			\mathcal{L}_j=&-F''({\bf{A}})\left[Q_s^u\sum_{m=1}^nM_{m}({\bf{G}})(|g_m|^2)_{x_j}\right]^2\\   
			&+F'({\bf{A}})Q_s^u\left(\sum_{m, q=1}^{n}M_{mq}({\bf{G}})\Big[(|g_m|^2)_{u_j}(|g_q|^2)_{u_j}-(|g_m|^2)_{x_j}(|g_q|^2)_{x_j}\Big] \right)\\
			&+F'({\bf{A}})Q_s^u\left(\sum_{m=1}^nM_{m}({\bf{G}})\Big[(|g_m|^2)_{u_ju_j}-(|g_m|^2)_{x_jx_j}+\frac{z_m^2-1}{k}\sum_{b=1}^{k_m}\bar{g}_mg_{m,bb}+\frac{\bar{z}_m^2-1}{k}\sum_{b=1}^{k_m}g_m\bar{g}_{m,bb}\Big]\right).
		\end{align*}
		Recall $F'>0$. Next, we set
		\begin{align*}
			\Psi_j=&-\frac{F''({M(\bf{G}}))}{F'(M({\bf{G}}))}\left[\sum_{m=1}^nM_{m}({\bf{G}})(|g_m|^2)_{x_j}\right]^2\\   
			&+\sum_{m, q=1}^{n}M_{mq}({\bf{G}})\Big[(|g_m|^2)_{u_j}(|g_q|^2)_{u_j}-(|g_m|^2)_{x_j}(|g_q|^2)_{x_j}\Big] \\
			&+\sum_{m=1}^nM_{m}({\bf{G}})\Big[(|g_m|^2)_{u_ju_j}-(|g_m|^2)_{x_jx_j}+\frac{z_m^2-1}{k}\sum_{b=1}^{k_m}\bar{g}_mg_{m,bb}+\frac{\bar{z}_m^2-1}{k}\sum_{b=1}^{k_m}g_m\bar{g}_{m,bb}\Big].
		\end{align*}
		We notice that $\frac{\mathcal{L}_j}{F'({\bf{A}})}\geq Q_s^u\Psi_j$. Indeed, this follows from Jensen's inequality, namely,  we have
		\begin{align*}
			- \frac{F''({\bf{A}})}{F'({\bf{A}})}(Q_s^uY)^2 &\geq- Q_{s}^{u}\left( \frac{F''(M({\bf{G}}))}{F'(M({\bf{G}}))}Y^2\right)  \qquad \text{for any intgrable} \ Y,
		\end{align*}
		since the function 
		\begin{align*}
			(t,y)\mapsto \frac{F''(t)}{F'(t)}y^2
		\end{align*}
		is convex on $(0,\infty)\times \R$. 

        Thus we have 
        \begin{align*}
            &C'(s)=\frac{1}{2}Q_{1-s}^{x}\left(\sum_{j=1}^{k} \mathcal{L}_{j}\right) \geq \frac{1}{2}Q_{1-s}^{x}\left(F'({\bf A})Q_{s}^{u} \left(\sum_{j=1}^{k} \Psi_{j} \right)\right).
        \end{align*}
        
        We will show that the inequality $\sum_{j=1}^k \Psi_j\geq0$ is the same inequality as the local condition of Proposition \ref{equivalComplex}. To show this, we express the three quantities inside the square brackets in the definition of $\Psi_j$ in terms of $g_{m,b}$. We write
		\begin{align*}
			\mathcal{Q}_1:&=\sum_{j=1}^k\left[\sum_{m=1}^n M_{m}({\bf{G}})(|g_m|^2)_{x_j}\right]^2\\
			&=\sum_{j=1}^k\sum_{m, q=1}^{n}\left[M_{m}({\bf{G}})(|g_m|^2)_{x_j} \right] \left[M_{q}({\bf{G}})(|g_q|^2)_{x_j} \right].
		\end{align*}
	For $m = 1, \ldots, n$, let $w_m \in \mathbb{C}^{k_m}$ be the complex vector with coordinates $w_{m,b} = \bar{g}_m g_{m,b} \in \mathbb{C}$, where $1 \leq b \leq k_m$. Plug in the partial derivatives $(|g_m|^2)_{x_j}$ we get
	\begin{align*}
		\mathcal{Q}_1&=4\sum_{j=1}^k\sum_{m, q=1}^{n}\sum_{\substack{1\leq b_1\leq k_m \\ 1\leq b_2\leq k_q}} (A_{m,b_1}\cdot e_j) (A_{q,b_2}\cdot e_j)M_{m}({\bf{G}}) M_{q}({\bf{G}})\Re(z_m \bar{g}_m g_{m,b_1}) \Re(z_q \bar{g}_q g_{q,b_2})\\
		&=4\sum_{m, q=1}^{n}\sum_{\substack{1\leq b_1\leq k_m \\ 1\leq b_2\leq k_q}} A_{m,{b_1}} A_{q,{b_2}}^T M_{m}({\bf{G}}) M_{q}({\bf{G}})\Re(z_m \bar{g}_m g_{m,b_1}) \Re(z_q \bar{g}_q g_{q,b_2})\\
		&=4\sum_{m, q=1}^{n} M_{m}({\bf{G}})M_{q}({\bf{G}})  (\Re(z_mw_{m}))^T A_{m} A_{q}^T  \Re(z_qw_{q}).
	\end{align*}
		Similarly, we also have that
		\begin{align*}
			\mathcal{Q}_2&:=\sum_{j=1}^k\left[(|g_m|^2)_{u_j}(|g_q|^2)_{u_j}-(|g_m|^2)_{x_j}(|g_q|^2)_{x_j}\right]\\
			&=4 \sum_{\substack{1\leq b_1\leq k_m \\ 1\leq b_2\leq k_q}} A_{m,{b_1}} A_{q,{b_2}}^T \left[\Re(\bar{g}_m g_{m,b_1})  \Re(\bar{g}_q g_{q,b_2})  -  \Re(z_m \bar{g}_m g_{m,b_1})  \Re(z_q \bar{g}_q g_{q,b_2})  \right]\\
			&=4\left( (\Re(w_{m}))^TA_{m} A_{q}^T  \Re(w_{q})  -  (\Re(z_mw_{m}))^T A_{m} A_{q}^T  \Re(z_qw_{q}) \right) .
		\end{align*}
		For the last quantity
		\begin{align*}
		\mathcal{Q}_3:= \sum_{j=1}^k\Bigg[\Bigg.&\sum_{b_1, b_2=1}^{k_m}\left[(A_{m,b_1}\cdot e_j) (A_{m,b_2}\cdot e_j)-z_m^2(A_{m,b_1}\cdot e_j) (A_{m,b_2}\cdot e_j)g_{m,b_1b_2}\bar{g}_m\right]\\
			+&\sum_{b_1, b_2=1}^{k_m}\left[(A_{m,b_1}\cdot e_j) (A_{m,b_2}\cdot e_j)-\bar{z}_m^2(A_{m,b_1}\cdot e_j) (A_{m,b_2}\cdot e_j)\bar{g}_{m,b_1b_2}g_m\right]\\
			+&\frac{z_m^2-1}{k}\sum_{b=1}^{k_m}g_{m,bb}\bar{g}_m +\frac{\bar{z}_m^2-1}{k}\sum_{b=1}^{k_m}\bar{g}_{m,bb}g_m+2(|(g_m)_{u_j}|^2 -|(g_m)_{x_j}|^2)\Bigg.\Bigg],
		\end{align*}
		we see that  moving the sum over $j$ inside the brackets, the first four terms vanish, leaving only the last term, since $A_mA^T_m=Id$. Also,
		\begin{align*}
			\sum_{j=1}^k|(g_m)_{u_j}|^2&=\sum_{j=1}^k\sum_{b_1, b_2=1}^{k_m}\Re((A_{m,b_1}\cdot e_j)g_{m,b_1}\overline{(A_{m,b_2}\cdot e_j)g_{m,b_2}})\\
			&=\sum_{b_1, b_2=1}^{k_m} A_{m,b_1}A_{m,b_2}^T\Re(g_{m,b_1}\bar{g}_{m,b_2})\\
			&=\sum_{b=1}^{k_m}|g_{m,b}|^2,
		\end{align*}
		and since $(g_m)_{x_j}=z_m(g_m)_{u_j}$, the quantity $	\mathcal{Q}_3$ simplifies to 
		\begin{align*}
			\mathcal{Q}_3=2\left(1-|z_m|^2 \right)\sum_{b=1}^{k_m}|g_{m,b}|^2.
		\end{align*}

        Thus we have 
        	\begin{align*}		\sum_{j=1}^k\Psi_j=&-4\frac{F''(M({\bf{G}}))}{F'(M({\bf{G}}))}\left(\sum_{1\leq m,q\leq n}  M_{m}({\bf{G}})M_{q}({\bf{G}}) (\Re(z_mw_{m}))^T A_{m} A_{q}^T  \Re(z_qw_{q})  \right)\\
			&+4\left(\sum_{ m, q=1}^{n} M_{mq}({\bf{G}})\left( (\Re(w_{m}))^T A_{m} A_{q}^T  \Re(w_{q})  -  (\Re(z_mw_{m}))^T A_{m} A_{q}^T  \Re(z_qw_{q}) \right)\right)\\
			&+2\left(\sum_{m=1}^n M_{m}({\bf{G}})\left(1-|z_m|^2 \right)\sum_{b=1}^{k_m}|g_{m,b}|^2 \right).
		\end{align*}
        Therefore the nonnegativity of $\sum_{j=1}^{j}\Psi_{j}$ follows from (\ref{Localcomplex}). Indeed, choose $c_{j} = |g_{j}|^{2}$, $j=1,\ldots, n$,  and $w_{m,b} = \bar{g}_{m}g_{m,b} \in \mathbb{C}$, $b=1, \ldots, k_{m}$, $m=1, \ldots, n$ in (\ref{Localcomplex}). 
	So we obtain $C'(s) \geq 0$. 


\eqref{monotonicitycomplex}$\implies$\eqref{Globalcomplex}:	Let $f_j: \mathbb{R}^{k_j} \to \mathbb{C}$, $j = 1, \ldots, n$, be polynomials expressed in the Fourier-Hermite basis, given by  
\begin{equation*}
	f_j(x) = \sum_{|\alpha| \leq \deg(f_j)} c_{\alpha,j} H_\alpha(x),
\end{equation*}
where $H_\alpha(x)$ denotes the Hermite polynomial indexed by the multi-index $\alpha$.  Now, consider the polynomials $\ell_j: \mathbb{R}^{k_j} \to \mathbb{C}$, $j = 1, \ldots, n$, expressed in the monomial basis as  
\begin{equation*}
	\ell_j(x) = \sum_{|\alpha| \leq \deg(f_j)} c_{\alpha,j} x^\alpha,
\end{equation*}
where $x^\alpha = x_1^{\alpha_1} x_2^{\alpha_2} \cdots x_{k_j}^{\alpha_{k_j}}$.
We observe, via a change of variables, that the function $C:[0,1]\to\R$ in \eqref{4Hflows} can be equivalently written as
\begin{equation*}
	C(s)=\int_{\R^k}F\left(\int_{\R^k}M(|q_1(u,x,s)|^2,\ldots,|q_n(u,x,s)|^2)\,d\gamma(u)\right)\,d\gamma(x),
\end{equation*}
where for $j=1,\ldots,n$,
\begin{equation*}
	q_j(u,x,s)=\int_{\R^{k_j}}\int_{\R^{k_j}}\ell_j(\sqrt{s}(A_j u+iv)+z_j\sqrt{1-s}(A_j x+iy))\,d\gamma(v)\,d\gamma(y).
\end{equation*}
Notice that
\begin{align*}
	q_j(u,x,0)&=T_{z_j}f_j(A_j x),\\
	q_j(u,x,1)&=f_j(A_j u).
\end{align*}
Hence, the function $C$ interpolates the inequality at the end points $s=0$ and $s=1$, namely
\begin{align*}
	C(0)&=	\int_{\R^k}F\circ M(|T_{z_1}(f_1(A_1 x))|^2, \ldots, |T_{z_n}(f_n( A_n x))|^2) \,d\gamma(x)\\
	C(1)&=F\left(	\int_{\R^k}M(|f_1(A_1 u)|^2, \ldots, |f_n(A_n u)|^2)\,d\gamma(u)\right),
\end{align*}
and therefore the monotonicity of $C$ finishes the proof.

	\end{proof}

\end{document}
